\chapter{Механизмы обучения}\label{ch:theory}

\section{Прямое распространение}

Сеть представляет собой цепочку полносвязных слоёв, разделённых нелинейностями. Для слоя
с входом $X$ размера $[N \times d_{\text{in}}]$, где $N$~--- число примеров в обрабатываемой
порции данных, выход вычисляется как
\begin{equation}
    Z = X W + b, \qquad A = f(Z),
\end{equation}
где $W$~--- матрица весов размера $[d_{\text{in}} \times d_{\text{out}}]$, $b$~--- вектор
смещений размера $[1 \times d_{\text{out}}]$, прибавляемый ко всем $N$ строкам, а $f$~---
функция активации, применяемая поэлементно.

На выходном слое применяется функция \texttt{softmax}, преобразующая вектор произвольных
вещественных чисел в распределение вероятностей по классам:
\begin{equation}
    p_i = \frac{\exp\bigl(z_i - \max_k z_k\bigr)}{\sum_j \exp\bigl(z_j - \max_k z_k\bigr)}.
    \label{eq:softmax}
\end{equation}
Вычитание максимума по строке не меняет результата, поскольку числитель и знаменатель
делятся на одну и ту же константу, однако предотвращает переполнение при вычислении
экспоненты от больших аргументов.

\section{Функция потерь}

В качестве функции потерь используется бинарная кросс-энтропия, вычисляемая по
вероятности положительного класса:
\begin{equation}
    L = -\frac{1}{N}\sum_{n=1}^{N}
        \Bigl[\, y_n \log p_n + (1 - y_n)\log(1 - p_n) \,\Bigr].
    \label{eq:bce}
\end{equation}
Вероятности предварительно ограничиваются отрезком $[\varepsilon,\ 1-\varepsilon]$ при
$\varepsilon = 10^{-12}$: в противном случае уверенный и при этом неверный прогноз даёт
$\log 0 = -\infty$, что приводит к появлению неопределённых значений во всех градиентах.
Помимо бинарной кросс-энтропии реализованы категориальная кросс-энтропия и
среднеквадратичная ошибка.

Деление на число примеров $N$ выполняется внутри функции потерь.
Благодаря этому величина шага градиентного спуска сохраняет смысл при любом размере
обрабатываемой порции данных.

\section{Обратное распространение через автоматическое дифференцирование}

Формулы обратного прохода для слоя в программе не записаны. Вместо них заданы правила
дифференцирования отдельных элементарных операций, а производные по всем параметрам
выводятся программно.

Каждая операция над тензором вычисляет не только результат, но и создаёт узел графа
вычислений, хранящий ссылки на узлы-аргументы и правило передачи градиента к ним. Граф,
таким образом, является побочным продуктом прямого прохода: отдельного этапа построения
графа не выполняется. Рёбра графа направлены от результата к аргументам, владение узлами
организовано в том же направлении, поэтому граф существует до тех пор, пока существует
ссылка на его корень, и освобождается целиком при её удалении.

Обратный проход состоит из трёх этапов. Сначала граф обходится в обратном топологическом
порядке; такой порядок гарантирует, что к моменту применения правила очередного узла в
его градиент уже просуммированы вклады всех узлов, которые этот узел использовали. Затем
в корень графа записывается начальное значение $\partial L / \partial L = 1$. После этого
каждое правило получает накопленный градиент и прибавляет свой вклад к градиентам
аргументов.

Накопление градиента выполняется суммированием. Поэтому узел,
использованный в нескольких ветвях сети, автоматически получает сумму вкладов, как того
требует правило дифференцирования сложной функции. Из этого же свойства следует, что
модель не обязана быть линейной цепочкой: на нём основана работа блока с обходным
соединением. Прямым следствием такого устройства является то, что базовый интерфейс
обучаемого блока содержит только прямой проход --- метод обратного прохода отсутствует у
всех слоёв.

Для слоя, задаваемого выражением $Z = XW + b$, описанный механизм приводит к следующим
выражениям:
\begin{equation}
    \frac{\partial L}{\partial W} = X^{\mathsf{T}}\,\frac{\partial L}{\partial Z},
    \qquad
    \frac{\partial L}{\partial b} = \sum_{n=1}^{N} \frac{\partial L}{\partial Z_{n,:}},
    \qquad
    \frac{\partial L}{\partial X} = \frac{\partial L}{\partial Z}\,W^{\mathsf{T}}.
    \label{eq:dense-grad}
\end{equation}
Градиент по вектору смещений представляет собой сумму по строкам, поскольку один и тот же
вектор $b$ участвует в вычислении всех $N$ строк результата.

Для поэлементных функций активации градиент умножается на производную соответствующей
функции:
\begin{equation}
    \sigma'(z) = \sigma(z)\bigl(1 - \sigma(z)\bigr), \qquad
    \tanh'(z) = 1 - \tanh^2(z), \qquad
    \mathrm{ReLU}'(z) = [\,z > 0\,].
    \label{eq:activation-grad}
\end{equation}

\section{Градиентный спуск}

После обработки каждой порции данных параметры смещаются в направлении, противоположном
градиенту:
\begin{equation}
    W \leftarrow W - \eta\,\frac{\partial L}{\partial W},
    \label{eq:sgd}
\end{equation}
где $\eta$~--- скорость обучения. Реализован также вариант с накоплением импульса,
сглаживающий траекторию спуска:
\begin{equation}
    v \leftarrow \mu\,v - \eta\,\frac{\partial L}{\partial W},
    \qquad
    W \leftarrow W + v.
    \label{eq:momentum}
\end{equation}

\section{Подготовка данных и организация обучения}

Набор данных содержит 569 примеров, по 30 вещественных признаков в каждом, при
соотношении классов 357 к 212 (62,74\,\% и 37,26\,\%). Признаки измеряются в
несопоставимых единицах: среднее квадратичное отклонение по столбцам различается более
чем на пять порядков. Поскольку градиент по весу пропорционален соответствующему входу,
без приведения признаков к общему масштабу подобрать единую скорость обучения для всей
сети невозможно. Поэтому признаки стандартизуются:
\begin{equation}
    x' = \frac{x - \mu}{\sigma},
    \label{eq:standardize}
\end{equation}
причём параметры $\mu$ и $\sigma$ вычисляются только по обучающей выборке и сохраняются
вместе с моделью. 

Выборка делится в соотношении 80 к 20 со стратификацией по классам, что даёт 456 примеров
для обучения и 113 для проверки. 

Одна эпоха представляет собой полный проход по обучающей выборке порциями по восемь
примеров. Перед каждой эпохой порядок примеров заново перемешивается, иначе
последовательность шагов спуска повторялась бы от эпохи к эпохе. Шаг обучения состоит из
обнуления градиентов, прямого прохода, вычисления функции потерь, обратного прохода и шага
оптимизатора. Обнуление необходимо именно потому, что градиенты накапливаются
суммированием. По завершении каждой эпохи вычисляются функция потерь и доля правильных
ответов на обеих выборках; эти вычисления выполняются в режиме, при котором граф не
строится. Дополнительно реализована ранняя остановка с возвратом к параметрам лучшей
эпохи.

Единое seed генератора случайных чисел используется при инициализации весов, разделении
выборки и перемешивании перед каждой эпохой, что обеспечивает воспроизводимость
результата. 
